\chapter{Corporate Stocks}

These companion notes follow Robert J. Shiller's Yale lecture on corporate stocks, with curation by LazyingArt LLC. We begin where the lecture begins, not with a formula but with a small company and a problem of fairness: who gets how many shares when one founder contributes money, another contributes time, and no one can value these contributions exactly in advance? From there the argument widens to the world scale of listed equity, then narrows again to boards, dividends, financing, and finally to Xerox and Microsoft, where the arithmetic of market capitalization and shareholder equity becomes concrete.

\section{Corporate Stocks as an Organizational Form}

We begin with the familiar picture of a corporation: people set up a company, divide it into shares, later issue more shares as the company grows, and sometimes sell those shares to someone else. The lecture reminds us that this form is old. The Dutch East India Company appears again as the first celebrated corporation with traded shares. The point is not antiquarian. The corporation is a durable scheme for organizing business.

Shiller then grounds that abstraction in his own experience. He helped found Case-Shiller Weiss, Incorporated, with Alan Weiss, Chip Case, and a businessman, Chuck Longfield. Only one of the four put money into the company; the others put in time and expertise. The immediate question was the one that always lurks behind the word \emph{equity}: how should the firm be divided? In the story, they did not pretend to solve that valuation problem exactly. They divided it four ways, partly out of fairness, partly because they wanted to begin.

That anecdote is not ornamental. It tells us what the corporate form is doing. A corporation lets heterogeneous contributions be bundled into a common project. It lets people stay involved at different intensities. It also gives the firm a way to manage resentment. Shiller says that when someone seemed to be doing more work than someone else, the natural salve was more shares, that is, a bonus. A share is therefore not just a claim on money at the end. It is also an instrument for motivating labor and settling conflict inside the organization.

This is why the lawyers matter so much in the story. Shiller says that founding a company gave him a new appreciation for lawyers because the charter anticipated conflicts before the founders fully understood what those conflicts would be. The charter and bylaws do not create the energy of the enterprise, but they keep the enterprise from dissolving as soon as the first resentments appear. That is the lecture's first deep point: the corporation is a social technology for combining money, labor, incentives, and conflict management.

He adds that the company was eventually sold to Fiserv in 2002, when it had only 12 employees, and that later the indexes themselves were acquired by S\&P and became widely known as the S\&P/Case-Shiller indexes. The company never became a giant firm, but that is not the lesson. The lesson is that even a small company already exhibits the essential corporate logic that we later see again, at much larger scale, in Microsoft or Xerox.

\section{The Scale of Listed Equity in the World Economy}

From the small company the lecture suddenly turns to the scale of listed equity around the world. The World Bank slide is important because it shows two different magnitudes at once: the dollar value of listed stocks and the value of those stocks relative to GDP.

\begin{figure}[h]
\centering
\includegraphics[width=0.9\textwidth]{lecture_09_figure_02.png}
\caption{World Bank stock-market capitalization table.}
\end{figure}

The cleaned mathematical object behind the slide is
\begin{equation}
\frac{\text{market capitalization}}{\text{GDP}}.
\end{equation}
The lecture immediately warns us not to read this too literally. We are comparing a stock with a flow: a stock of wealth in the stock market divided by one year's GDP. The ratio is still informative, but only if we remember what is being divided by what.

A few visible rows from the slide are enough to carry the argument:
\[
\begin{array}{lcc}
\text{Region} & \text{Fraction of GDP} & \text{US\$ trillions} \\
\text{United States} & 0.817 & 11.737 \\
\text{United Kingdom} & 0.696 & 1.851 \\
\text{Canada} & 0.669 & 1.002 \\
\text{East Asia and Pacific} & 0.559 & 3.244 \\
\text{Euro Area} & 0.380 & 5.153
\end{array}
\]
The lecture's emphasis is that the United States is not merely large in absolute terms. It is also large relative to GDP. In 2008 its listed market capitalization was worth almost a full year's GDP. The United Kingdom is also large relative to GDP, while the Euro Area is large in dollar magnitude but much smaller on this stock-to-flow comparison.

Shiller connects that scale to history. The United States became a major capitalist center early, and he points to the New York corporate law of 1811 as part of the legal background that helped set the form in motion. He also expects the chart eventually to look different. Stock markets are spreading and growing elsewhere because, in his view, the corporate form works. Even countries that once defined themselves against capitalism have adopted stock markets in some form. The lecture uses that historical spread as evidence that the institutional device has been learned and retained.

But almost immediately Shiller pulls back. He does not want to exaggerate the size of the stock market. If we round U.S. listed market value to \$12 trillion and divide by roughly 300 million people in 2008, then
\begin{equation}
\frac{12\ \text{trillion dollars}}{300\ \text{million people}}
\approx \$40{,}000 \text{ per person}.
\end{equation}
For a family of four this becomes about
\begin{equation}
4 \times \$40{,}000 = \$160{,}000.
\end{equation}
That is not nothing, but it is not the whole economy. Shiller presses the point further by comparing stocks with housing wealth. He says that U.S. houses were worth about \$18 trillion in 2010, more than the stock market, and that total household assets were about \$69 trillion. Stocks are important, but they are only one claim on national wealth.

The lecture also insists on why the stock market is smaller than many people vaguely imagine. Corporate stocks are only claims on profits after many other payments have already been made. Corporations pay wages, interest, taxes, and many other expenses before common shareholders receive what remains. So even in a highly capitalist economy, the stock market is not the whole social product. It is the capitalized value of a residual claim.

\section{Corporation, Equity, and Governance}

At this point the lecture narrows again and asks what sort of thing a corporation actually is. Before that, though, it makes a useful distinction. The World Bank slide is about traded stocks. Not every corporation is publicly traded. Shiller's own company was never traded; that is private equity. Public equity means shares traded on an exchange. He gives the German terminology as a reminder that this distinction is not peculiar to the United States: \emph{AG}, \emph{Aktiengesellschaft}, indicates a stock-market company, while \emph{GmbH} indicates the nontraded form.

He then pauses over the word \emph{equity}. In ordinary English the word has many meanings, but here it means shares in a corporation. He speculates that the term evokes equality or fairness, and institutionally that is a useful guide. Common shares are, in the ordinary case, treated equally. That is the financial meaning that matters in the lecture.

The next term is \emph{corporation} itself. From the Latin \emph{corpus}, body, the corporation is a legal person. Not a natural person, but an artificial one. It can own property, incur liabilities, make contracts, and continue over time even as the human beings around it enter or exit. Shiller even notes that corporations in some form go back to Rome, where \emph{publicani} had something like a stock market in the Roman Forum, though corporate proliferation on a mass scale is mainly a nineteenth-century phenomenon.

Once we have legal personhood, governance follows. The standard picture is simple:
\begin{itemize}
\item the shareholders own the corporation through shares;
\item each share normally carries one vote;
\item shareholders elect a board of directors;
\item the board hires and supervises the chief executive officer.
\end{itemize}
This is what the lecture calls shareholder democracy. It is modeled on political democracy, but it is not one-person-one-vote. It is one-share-one-vote. That difference matters, because control is tied to ownership.

Shiller makes the institutional details concrete. In the United States one chooses a state in which to incorporate. Delaware is popular because of its corporate law and tax environment. State law typically requires an annual shareholder meeting. Yet the shareholders do not run the company day to day. The board does. The board is small, perhaps six to ten people, and is expected to know the company well enough to judge management. Directors are not generally the employees who run the firm. They are the people brought in to supply judgment, perspective, and oversight. The CEO, by contrast, is the energetic full-time executive, working long hours and reporting to the board.

The board is not meant to be decorative. It can replace the CEO. That is why the lecture brings in Carl Icahn. The point is not to celebrate every activist investor, but to show that shareholder control can become real when a large shareholder pushes a complacent board to act. Boards owe duties of loyalty to shareholders, not merely friendliness toward management.

The lecture then turns to the contrast between for-profit and nonprofit corporations. Yale itself is the nearby example. Yale is a nonprofit corporation: it has a board, legal personhood, and an institutional mission, but it has no shareholders entitled to profits. Any surplus stays inside the institution and is used for education and research. That sounds very different from a for-profit corporation, yet Shiller emphasizes that the organizational sociology can be surprisingly similar. The same kinds of people can sit on both sorts of boards, and they may behave in similar ways. The distinction is not that one has structure and the other does not. The distinction is that in the for-profit case there are owners with claims on distributed profits.

\section{Shares, Market Capitalization, and the Logic of Dividends}

Only after all this structure is in place does the lecture ask the direct financial question: what is a share? The arithmetic is elementary, but Shiller deliberately repeats it, because students and investors alike often lose sight of the denominator.

Let $P_t$ be the price per share at time $t$, $N_t$ the number of shares outstanding, and $\mathrm{MC}_t$ the market capitalization. Then
\begin{equation}
\mathrm{MC}_t = P_t N_t.
\end{equation}
Equivalently,
\begin{equation}
P_t = \frac{\mathrm{MC}_t}{N_t}.
\end{equation}
If we own $n_t$ shares, our ownership fraction is
\begin{equation}
\frac{n_t}{N_t}.
\end{equation}

This is why the raw number of shares owned is meaningless without the total number outstanding. If a company has 1,000,000 shares and we own 1,000, then
\[
\frac{1000}{1{,}000{,}000} = \frac{1}{1000}.
\]
If instead the company has 10,000,000 shares, then the same 1,000 shares amount only to
\[
\frac{1000}{10{,}000{,}000} = \frac{1}{10{,}000}.
\]
That corrected arithmetic is essential, because the lecture keeps returning to ownership fractions, not raw share counts.

The lecture also insists that common shares are treated equally. If the firm pays anything to one common shareholder, it must pay the other common shareholders in proportion to their holdings. That is the practical meaning of equal shares.

This leads naturally to dividends. A dividend is paid only when the board decides to pay it. Shiller remembers exactly such a conversation in his own small company. Employees had been given tiny amounts of stock as motivation, and the board wondered whether it should pay a dividend simply so that those employees would receive checks in the mail and feel that their shares were worth something. That anecdote matters because it places us exactly where the lecture wants us: between the legal discretion of the board and the economic point of owning stock.

\subsection{Question \& Answer}

\paragraph{Question.} If boards are not legally required to pay dividends, why do common shares have value at all?

\paragraph{Answer.} The lecture's answer is that the whole point of a for-profit corporation is eventually to distribute profits to its owners. A young company may postpone that day. Microsoft is the example Shiller keeps in reserve: for a long time it paid no dividend and simply kept the money, because shareholders believed that the real opportunities still lay ahead. But a company that would never distribute anything would cease, in economic substance, to look like a for-profit company at all. Shareholders tolerate delay, not permanent denial.

That is why the lecture corrects a common misunderstanding. People say they buy stocks because the price will go up. Shiller says that in theory this is backward. The price should go up only because investors expect more dividends later. The benchmark statement is
\begin{equation}
P_t = \sum_{j \ge 1}\frac{\mathbb{E}_t[\mathrm{Div}_{t+j}]}{(1+r)^j}.
\end{equation}
The lecture does not derive this formula formally. It uses it as the disciplinary idea behind the rest of the discussion. A share is valuable because it is a claim on expected future payouts.

The remaining worry is board discretion. What if the board simply hoards the cash? Shiller's answer is governance. Shareholders elect the board. If the board accumulates money and refuses to distribute it despite strong profits, shareholders can replace the board. In the lecture, Carl Icahn reappears here in thought if not in person: put someone like that on the board and the money starts to come out. So the present-value logic is not detached from institutional control. It depends on it.

\section{Dividend Pricing and Financing Mechanics}

Once that puzzle has been resolved, the next piece is almost embarrassingly simple. If the company pays cash out, the company is worth less by that amount. This is the logic of the ex-dividend price.

Let the aggregate cash dividend at time $t$ be $D_t$, with $N_t$ shares outstanding. Then the dividend per share is
\begin{equation}
d_t = \frac{D_t}{N_t}.
\end{equation}
The lecture's pricing logic can be written as a short derivation:
\begin{enumerate}
\item before the dividend, the firm owns the cash that is about to be distributed;
\item the board pays out aggregate cash $D_t$;
\item assets fall by $D_t$, while liabilities do not automatically change;
\item therefore equity value falls by roughly $D_t$;
\item per share, the price should fall by roughly $d_t$.
\end{enumerate}
So the ex-dividend adjustment is
\begin{equation}
P_t^{\mathrm{ex}} \approx P_t^{\mathrm{cum}} - d_t.
\end{equation}

Shiller treats this as everyday arithmetic, not a deep theorem. In the old newspapers, and now on websites, the ex-dividend adjustment is marked because the observed price decline is not bad news. It is merely cash leaving the company.

This is also why ``selling dividends'' is misleading and, in the lecture's telling, ethically dubious. A broker should not urge a client to rush into a stock merely because a dividend will be paid in three days. If one buys before the dividend, one gets the dividend but pays the cum-dividend price. If one buys after the dividend, one misses the dividend but buys the stock more cheaply. Apart from taxes and second-order frictions, it comes to about the same thing.

From there the lecture moves back to the small-company picture. Shares are valuable partly because they are saleable. If one founder gets tired and wants out, the firm or the other founders can buy back the shares. If a new investor wants in, the firm can sell shares at the current price, which may be very different from the price at which the company was first organized. Saleability lets the organization persist even when particular people come and go.

The same logic explains financing. If the company needs a new factory, it can issue more shares and raise money. In a private company there may be no exchange quote, but a broker can still search for wealthy buyers or other investors willing to purchase the new issue. This is equity financing. The alternative is debt financing, by bank loan or bond issue. Debt carries fixed repayment obligations; equity carries a residual claim on profits. The lecture plainly intends that contrast, despite a garbled line in the transcript.

Once the firm borrows, common shareholders become more exposed. Their claim is the residual
\begin{equation}
\mathrm{SE}_t = A_t - L_t,
\end{equation}
where $A_t$ denotes assets and $L_t$ liabilities. Equity is what remains after creditors have been paid. This is the lecture's leverage point: debt does not merely add financing, it also amplifies the risk borne by common shareholders because their slice is now thinner and comes later in the priority structure.

\subsection{Question \& Answer}

\paragraph{Question.} Why is dilution not automatically bad, and why is a buyback not automatically good?

\paragraph{Answer.} Because neither event can be judged from the share count alone. Suppose a firm has 1,000,000 shares outstanding and we own 1,000 of them. Then we own one-thousandth of the firm. Now suppose the board issues another 1,000,000 shares at \$10 per share. The lecture's arithmetic is
\begin{align}
N &= 1{,}000{,}000, & n &= 1000, & \frac{n}{N} &= \frac{1}{1000}, \\
N' &= 2{,}000{,}000, & n &= 1000, & \frac{n}{N'} &= \frac{1}{2000}.
\end{align}
Our ownership fraction has been cut in half. That is dilution. But the firm has also received
\begin{equation}
1{,}000{,}000 \times \$10 = \$10{,}000{,}000
\end{equation}
in cash. So the right question is not simply whether our fraction fell. The right question is whether the firm sold the new shares at a good enough price and whether the cash will be used well.

The reverse side is repurchase. If the firm buys back shares, the share count falls. If it buys back 5\% of all shares and we tender 5\% of our own holdings, then we end up with cash in hand and essentially the same ownership fraction in what remains. That is why Shiller says that a buyback can feel much like a dividend. One historical reason firms preferred repurchases was tax treatment: dividends were taxed as income while capital gains could be taxed more lightly. By the time of the lecture, those rates were closer together, but tax considerations still remained complicated enough to matter.

The lecture then uses this financing discussion to introduce common and preferred stock. Common stock is the real ownership claim, the ultimate residual claimant. Preferred stock sits above common. It usually promises a fixed dividend, often has no ordinary vote, and typically must be paid up before common dividends resume. It is therefore safer than common stock but lacks the full upside of common. The government purchase of preferred shares in General Motors during the bailout illustrates the point. Preferred shares let the government provide capital without looking as though it had fully nationalized the company by taking ordinary voting control.

\section{Do Stock Markets Matter for Big Firms, and How Do Firms Choose Dividends?}

At this point the lecture tests its own story. It is easy to see how a small company of founders, employees, and early investors uses shares to organize itself. But what about a giant company? Bill Gates and Paul Allen started Microsoft in 1975 as a little company. Once a firm grows very large, does the stock market still play the role we have been describing, or does it become mostly a market in old shares?

Shiller raises the challenge in several forms. He briefly invokes the old Marxian suspicion that stock markets may become detached from real production. He then turns to the modern finance version of the same complaint, Stuart Meyers's pecking-order theory. On Meyers's aggregate U.S. data for 1973--1982, the financing mix looked roughly like
\begin{equation}
\text{financing mix} \approx 62\%\ \text{retained earnings} + 6\%\ \text{net equity} + \text{remainder debt}.
\end{equation}
That picture suggests a hierarchy. Firms prefer to use their own retained earnings first, then debt, and only as a last resort issue new shares. If that were the whole story, one might conclude that large firms hardly use the stock market for financing at all.

\subsection{Question \& Answer}

\paragraph{Question.} If big firms often finance internally, what is the stock market still doing for capitalism?

\paragraph{Answer.} The lecture's answer is that the aggregate net number is too crude. What Meyers calls 6\% equity financing is \emph{net} equity issuance:
\begin{equation}
\text{net equity issuance}
=
\text{new shares sold}
-
\text{share repurchases}.
\end{equation}
A low net number can arise even when many firms are issuing shares, provided many others are buying them back.

That is exactly the Fama--French reply. They argue that Meyers was looking at aggregate data during a period of relatively low net issuance. But when one asks how many firms actually issue equity in a given year, the picture changes. Shiller quotes their result that 67\% of U.S. firms issued new shares each year during 1973--1982, and 86\% did so during 1993--2002. So the stock market still matters. It matters for financing, and it matters for price discovery even when mature firms are not constantly flooding the market with new stock.

Once that question is settled, the lecture turns back from market-wide challenge to boardroom behavior. Even if firms retain earnings, borrow, issue, and repurchase in different mixtures, someone still has to decide what dividend to pay. That is where John Lintner enters. The route is important: the mathematics comes only after the institutional behavior has been described.

Lintner, a Harvard Business School professor and one of the early figures behind the capital asset pricing model, interviewed managers and directors about dividend decisions. He found a great deal of verbal complexity. Boards want the company to be well regarded. Mature firms take pride in paying dividends. Above all, boards dislike cutting dividends once they have started. That single behavioral regularity does a great deal of work in the lecture.

\begin{figure}[h]
\centering
\includegraphics[width=0.78\textwidth]{lecture_09_figure_03.png}
\caption{Lintner dividend adjustment on the board.}
\end{figure}

The board image is valuable precisely because it catches the equation in mid-emergence. On the left the earlier labels ``common,'' ``preferred,'' and ``John Lintner'' are still visible from the surrounding discussion. In the center we can clearly read $\mathrm{Div}_t-\mathrm{Div}_{t-1}$ and the beginning of the right-hand side, $=\rho$. The full equation is not completely legible on the board, so the displayed formula below is a cautious transcript-assisted reconstruction:
\begin{equation}
\mathrm{Div}_t - \mathrm{Div}_{t-1}
=
\rho\bigl(\tau\,\mathrm{EPS}_t - \mathrm{Div}_{t-1}\bigr)
+
\varepsilon_t,
\end{equation}
with
\[
0 < \tau < 1,
\qquad
0 < \rho < 1.
\]

The interpretation is simple and powerful:
\begin{enumerate}
\item the board has a target payout ratio, so the desired dividend level is roughly $\tau\,\mathrm{EPS}_t$;
\item the board dislikes cuts, so it does not jump all the way to target when earnings rise;
\item instead, dividends move only part of the way toward target, at speed $\rho$;
\item the disturbance term $\varepsilon_t$ absorbs the boardroom complexity that the model does not explicitly represent.
\end{enumerate}

This is the lecture's clearest mathematical core. The equation does not drop from the sky. It comes from an institutional fact: firms fear the reputational cost of cutting dividends. If earnings surge, the board raises dividends cautiously, because an aggressive increase today may force an embarrassing cut tomorrow. Shiller also notes that payout ratios have trended downward over long stretches of time. Earlier shareholders wanted a larger immediate fraction of earnings. Modern shareholders are often more willing to let firms retain cash for internal growth.

\section{Xerox and Microsoft: Book Equity and Market Value}

At this point the lecture does what it promised earlier and turns to concrete firms. Xerox comes first. The story is almost mythic: founded in 1906 as the Halloyd Photographic Company, transformed by Chester Carlson's invention of the plain-paper dry copier, and elevated into one of the glamour growth firms of the 1960s and 1970s. Then came digital copying, which Xerox did not own, and the firm nearly collapsed. Ann Mulcahy, hired in 2001, is presented as the executive who kept the company from extinction.

Because Xerox is a public company, its balance sheet is visible in SEC filings. Shiller uses that institutional detail to make a larger point. Going public gives a firm access to public markets, but it also subjects it to continuous disclosure. His own private company never had to publish a quarterly balance sheet on \texttt{sec.gov}. Xerox did.

The balance-sheet lesson is basic but essential. Assets are what the firm owns. Liabilities are what it owes. Receivables, inventory, buildings, short-term debt, long-term debt, deferred taxes, and even a small amount of preferred stock all appear in the lecture's summary. What finally matters for common shareholders is the residual, the shareholder's equity, also called stockholder's equity. In 2010 Xerox's shareholder equity was about \$11.9 billion, up from roughly \$4.9 billion. That is the book residual after liabilities are subtracted from assets.

But Shiller is careful. Shareholder equity is an estimate, not a market price. Buildings and equipment are carried at accounting values, and no one knows exactly what they would fetch if sold off in a hurry. So book equity is a liquidation-style benchmark, not the same thing as market value.

For Xerox the market-cap arithmetic is
\begin{align}
\text{Xerox market cap}
&\approx (1.4\ \text{billion shares})(\$11.80) \\
&\approx \$16.4\text{--}16.5\ \text{billion}.
\end{align}
So the market valued Xerox above its book shareholder equity of about \$11.9 billion.

Microsoft sharpens the contrast. Founded in 1975 by two young founders, it had by 2010 assets of about \$92 billion, including a great deal of cash and short-term investment, some real estate and equipment in Redmond, and even about \$10 billion of stock holdings. On the liability side Shiller mentions ordinary items such as accounts payable, taxes, and unearned revenue, the latter being money already received on products for which future service is still owed. After subtracting liabilities, Microsoft had shareholder equity of about \$48 billion.

Yet the market capitalization was far larger:
\begin{align}
\text{Microsoft market cap}
&\approx (8.497\ \text{billion shares})(\$26) \\
&\approx \$221\ \text{billion}.
\end{align}
That is more than five times book shareholder equity. Shiller even suggests that Microsoft lends some support to the pecking-order intuition: for a long time it retained cash and paid no dividend, and investors were willing to accept that because they believed in the firm's opportunities and management.

\subsection{Question \& Answer}

\paragraph{Question.} Why can market capitalization exceed shareholder equity on the books?

\paragraph{Answer.} Because a public company is not purchased merely as scrap. If we buy the company in order to shut it down, sell the buildings, collect the receivables, unload the inventory, and pay off the debts, then book shareholder equity is the rough benchmark for what might be left. But if we buy the company as an ongoing concern, then we are buying more than booked assets. We are buying future profits, organizational capital, reputation, installed customer relationships, and growth opportunities.

That is why Xerox can sell above its book equity even after severe competitive trouble. The market is saying the company is worth more alive than dead. And that is why Microsoft can sell vastly above its book equity. Its value lies overwhelmingly in future operating profits, not in liquidation of the Redmond campus and short-term securities alone.

The lecture also points out the converse. If market capitalization fell far below shareholder equity, the firm could become vulnerable to takeover and liquidation. Buy the company cheaply, shut it down, sell the assets, pay the debts, and keep the difference. That is why CEOs and boards watch market capitalization so closely against book equity. Still, the comparison is never exact, because book shareholder equity is itself an accounting estimate. It is informative, but it is not a literal auction value observed every day.

\section{Summary}

The lecture unfolds in a very deliberate order. We begin with the corporation as an organizational invention, one that lets people combine money, labor, incentives, and conflict resolution inside a durable legal form. We then widen to the world scale of listed equity, but only long enough to see both its importance and its limits. After that we rebuild the institution from the inside: legal personhood, public versus private equity, one-share-one-vote, boards, CEOs, and the contrast between for-profit and nonprofit corporations.

Only then do we turn fully to valuation. Market capitalization is price times shares outstanding. Ownership is a fraction, $n_t/N_t$, not a raw share count. Common stock is ultimately a claim on dividends, even if the board delays those dividends for years. The ex-dividend price fall, dilution arithmetic, buybacks, leverage, and preferred stock all become simple once we keep the residual nature of common equity in view. The stock market still matters even for big firms, because net issuance is not the same thing as inactivity, and price discovery remains central. Lintner's model then distills messy boardroom judgment into a partial-adjustment rule for dividends. Xerox and Microsoft close the chapter by showing why book shareholder equity and market capitalization can differ so sharply. The corporation is not just a legal shell and not just a balance sheet. It is a going concern whose claims are priced through expectations about future distributions.